\documentclass[12pt]{article}
\usepackage[utf8]{inputenc}
\usepackage[T1]{fontenc}
\usepackage{hyperref}
\title{Uripiv-Wala-Rano-Atchin Language Resources for K-12 Curriculum Development}
\author{Compiled by Claude Sonnet 5 (claude-sonnet-5) for Andr\'as Kornai}
\date{}
\begin{document}
\maketitle

\section{Introduction}
Uripiv-Wala-Rano-Atchin (ISO 639-3: upv), also called Northeast Malakula, is an Austronesian dialect chain on the northeast coast of Malekula, Vanuatu. It is spoken on the islands of Uripiv, Wala, Rano and Atchin and on the opposite mainland. Estimates give about 9,000 speakers. Wikipedia reports literacy in the language at 10--30\%. This search found little in the way of educational materials.

\section{Primary Educational Resources}
No primers, textbooks or school programmes were found in the search results.

\section{Linguistic and Reference Materials}
\begin{itemize}
\item Wikipedia on Uripiv and on the Northeast Malakula language: \url{https://en.wikipedia.org/wiki/Northeast_Malakula_language} and \url{https://en.wikipedia.org/wiki/Uripiv}
\item Austronesian Basic Vocabulary Database word lists: \url{https://abvd.eva.mpg.de/austronesian/language.php?id=1146} (Wala) and \url{https://abvd.eva.mpg.de/austronesian/language.php?id=1069} (Uripiv)
\item Numerals page, Max Planck Institute: \url{https://lingweb.eva.mpg.de/channumerals/Uripiv-Wala-Rano-Atchin.htm}
\end{itemize}

\section{Religious and Community Resources}
\begin{itemize}
\item Joshua Project profile (Uripiv-Atchin): \url{https://joshuaproject.net/people_groups/15734}
\end{itemize}

\section{Digital and Audio Resources}
\begin{itemize}
\item Global Recordings Network page for the language: \url{https://globalrecordings.net/en/language/upv}
\item Global Recordings Network page for ``Malakula: East'': \url{https://globalrecordings.net/en/language/4714}
\end{itemize}

\section{Curriculum Development Context}
Little curriculum-ready material exists. The word lists and recordings are a starting point, but a native-language curriculum would need new orthography, reading and teaching materials, ideally developed with local speakers.
\end{document}